\section{Introduction}
Deploying an object detector in a new domain forces a practitioner to choose \emph{how} to spend an annotation budget. Three families of strategies are now realistic. (i)~An open-vocabulary detector such as YOLO-World~\cite{cheng2024yoloworld}, Grounding~DINO~\cite{liu2024groundingdino}, OWLv2~\cite{minderer2023owlv2} or YOLOE~\cite{wang2025yoloe} can be prompted with class names and used \emph{zero-shot}, at zero labeling cost. (ii)~Such a detector can act as a \emph{teacher} that pseudo-labels an unlabeled image pool, from which a small, fast \emph{student} (e.g., YOLO) is trained~\cite{son2024teacherstudent,xin2024dart}. (iii)~A pretrained closed-set detector can be \emph{fine-tuned} on manually labeled images. Which of these is best depends on the number of labeled images that can be afforded, on how far the target label space is from what the open-vocabulary detector has seen, and on the deployment constraints on latency.

The literature mostly improves one strategy at a time: stronger open-vocabulary architectures~\cite{cheng2024yoloworld,liu2024groundingdino,minderer2023owlv2,wang2025yoloe}, better pseudo-labeling~\cite{gao2022pseudobox,liu2021unbiased,xu2021softteacher}, or few-shot fine-tuning recipes~\cite{wang2020tfa}. Recent benchmarks show that zero-shot detectors can fail badly on concepts that are rare in pre-training data~\cite{robicheaux2025rf100vl}, and that pseudo-labels alone can rival supervised training in some domains~\cite{abid2025waste}. What is missing is a \emph{budget-indexed} comparison: given $N$ labeled images, which strategy should be used, how reliable is that decision, and can it be made \emph{before} the labeling money is spent?

This paper makes four contributions.
\begin{itemize}
\item \textbf{Budget-controlled benchmark.} We compare zero-shot detection, teacher--student pseudo-labeling, and fine-tuning under identical evaluation, nested labeled subsets ($N\in\{0,10,\dots,400\}$ images), matched training compute, and multiple seeds, on two domains that differ in how well the zero-shot teacher fits them, and report the budgets at which the ranking of strategies changes.
\item \textbf{Gold-calibrated pseudo-labeling (TS-cal).} We use the same $N$ labeled images that fine-tuning would consume to select the teacher's pseudo-label confidence threshold, so that labeled data serve both as supervision and as a calibration set for the teacher.
\item \textbf{Cost-aware strategy selector.} We fit per-strategy learning curves on small pilot budgets and extrapolate to choose a strategy for a larger budget, and we evaluate the choices against measured outcomes (regret relative to an oracle and to fixed rules).
\item \textbf{Pilot-set reliability.} We quantify how many labeled images are needed merely to \emph{measure} the zero-shot detector's accuracy and to order it against a fine-tuned model, a prerequisite for any budget-aware decision.
\end{itemize}
All code, raw per-image statistics, and scripts are released with the paper.

\section{Related Work}
\textbf{Open-vocabulary detection.} GLIP~\cite{li2022glip} casts detection as phrase grounding; OWL-ViT/OWLv2~\cite{minderer2022owlvit,minderer2023owlv2}, ViLD~\cite{gu2022vild} and Detic~\cite{zhou2022detic} transfer image-level vision--language knowledge (e.g., CLIP~\cite{radford2021clip}) to box prediction; Grounding~DINO~\cite{liu2024groundingdino} fuses language early in a DETR-style detector. YOLO-World~\cite{cheng2024yoloworld} and YOLOE~\cite{wang2025yoloe} bring open-vocabulary prompting to the real-time YOLO family. Our zero-shot detector is YOLOE with the MobileCLIP~\cite{vasu2024mobileclip} text encoder. Domain-specific adaptations of YOLO-World to retail scenes and one-shot retail detectors such as SiamYOLOv8~\cite{desmarescaux2025siamyolov8} target accuracy rather than the choice between strategies.

\textbf{Pseudo-labeling and teacher--student detection.} Semi-supervised detectors such as Unbiased Teacher~\cite{liu2021unbiased} and Soft Teacher~\cite{xu2021softteacher} bootstrap from a labeled subset. With vision--language teachers, Gao \emph{et al.}~\cite{gao2022pseudobox} generate pseudo boxes for open-vocabulary training; Son and Jung~\cite{son2024teacherstudent} train YOLO on Grounding~DINO labels and report no clear degradation relative to manual labels; DART~\cite{xin2024dart} automates a diversification--annotation--review--training pipeline; and Abid \emph{et al.}~\cite{abid2025waste} show that pseudo-labels from vision--language models can improve YOLO on waste detection. These works study one pipeline at fixed supervision; none varies the manual-label budget to locate where each strategy wins.

\textbf{Benchmarks of zero-shot versus few-shot detection.} Roboflow100-VL~\cite{robicheaux2025rf100vl} evaluates vision--language detectors on 100 diverse datasets in zero-, few-shot, semi-supervised and fully supervised regimes and shows large zero-shot failures outside common concepts. Our study is narrower but complementary: it sweeps the labeling budget densely within a domain, controls training compute, and turns the measurements into a selection rule.

\textbf{Annotation cost and learning curves.} Box annotation is costly: on the order of seconds per box even with efficient protocols~\cite{su2012crowdsourcing,papadopoulos2017extreme}. Active learning for detection~\cite{brust2018active} decides \emph{which} images to label; we address the complementary question of \emph{which learning strategy} to buy with a given number of labels. Our selector relies on the empirical regularity of learning curves~\cite{cortes1993curves,hestness2017scaling}.
